\documentclass[11pt]{article}

\usepackage{amsmath,amssymb}
\usepackage{xurl}
\usepackage[colorlinks=true,linkcolor=blue,citecolor=blue,urlcolor=blue]{hyperref}

\input{command}

\title{Source-General Projected Support Corner Star-Algebra}
\author{}
\date{2026-07-30}

\begin{document}
\maketitle
\gapnote{scope-restriction}{resolved}

\section{The assertion}

Let $T:\MN{D}\to\MN{D}$ be positive and trace preserving. Let $T^*$ denote
its trace adjoint, assume that $T^*$ is unital and satisfies the Schwarz
inequality, and let $\rho\succeq0$ be a maximum-rank fixed point of $T$. If
$Q$ is the support projection of $\rho$, Corollary~6.6 of
Ref.~\cite{Wolf2012Quantum} states that
\begin{align}
    \mathcal A_Q
        &=
        \left\{Y\in Q\MN{D}Q
            \mathrel{}\middle|\mathrel{}
            QT^*(Y)Q=Y\right\}
    \label{eq:wolf_cor66_kraus_hypothesis_restriction_corner_fixed_set}
\end{align}
is a $*$-subalgebra of the corner algebra $Q\MN{D}Q$. The local source is
\path{Notes/WolfNoteTexSource/ch06_spectral_properties.tex},
lines~1423--1437.

Complete positivity is an example of a sufficient hypothesis in
Ref.~\cite{Wolf2012Quantum}; it is not an assumption of Corollary~6.6. For a
unital map, complete positivity implies the Schwarz inequality, but the
Schwarz condition is strictly weaker because some positive maps that are not
completely positive also satisfy it.

\section{The former restriction}

The earlier formal result was
\leanid{Kraus.cornerFixedPointsStarSubalgebra}, with membership
characterization
\leanid{Kraus.mem\_cornerFixedPointsStarSubalgebra}, in
\path{QICLean/Channel/FixedPoint/CornerFixedPoints.lean}. Its hypotheses are
a Kraus family $K$ satisfying \leanid{IsTP} $K$ and a positive semidefinite
fixed point of the induced map.

A Kraus representation assumes complete positivity. The earlier declarations
therefore had strictly stronger hypotheses than Corollary~6.6 of
Ref.~\cite{Wolf2012Quantum}; they remain completely positive
specializations, not the source-facing formalization.

The analogous former restriction for Corollary~6.7 of
Ref.~\cite{Wolf2012Quantum} is recorded in
\path{wolf_cor67_maximal_support_restriction.tex}.

\section{Resolution}

The source-general result is
\leanid{IsPositiveMap.stationaryCornerAdjointFixedPointsStarSubalgebra}, with
membership characterization
\leanid{IsPositiveMap.mem\_stationaryCornerAdjointFixedPointsStarSubalgebra},
in
\path{QICLean/Channel/FixedPoint/AbstractCornerFixedPoints.lean}. Its
hypotheses are exactly positivity and trace preservation of $T$ together with
the Schwarz inequality for $T^*$. It assumes neither a Kraus representation
nor complete positivity. The result holds for every positive semidefinite
fixed point, so it includes the maximum-rank fixed point used in
Corollary~6.6 of Ref.~\cite{Wolf2012Quantum}.

The proof follows the compression argument in lines~1433--1437 of the local
source for Ref.~\cite{Wolf2012Quantum}. Compress $T$ to the support of
$\rho$. The compressed map is positive and trace preserving, has a positive
definite fixed point, and has a Schwarz trace adjoint. The abstract fixed-point
algebra theorem \leanid{SchwarzMap.fixedPointsStarSubalgebra} therefore
applies.

Extension through the support isometry transports multiplication from the
compressed algebra to $Q\MN{D}Q$. The trace-adjoint compression identity
identifies the compressed fixed-point equation with
\begin{align}
    QT^*(Y)Q
        &=Y.
    \label{eq:wolf_cor66_kraus_hypothesis_restriction_compressed_fixed_equation}
\end{align}
This proves that the set
(\ref{eq:wolf_cor66_kraus_hypothesis_restriction_corner_fixed_set}) is a
$*$-subalgebra under the source hypotheses.

\section{Conclusion}

The complete-positivity restriction for Corollary~6.6 of
Ref.~\cite{Wolf2012Quantum} is eliminated. The blueprint and numbering
coverage now point to the source-general declaration. The Kraus declarations
remain documented completely positive specializations.

\bibliographystyle{alpha}
\bibliography{references}

\end{document}
